\chapter{Introducción}
\label{cap:introduccion}

\section{Contexto y motivación}
\label{sec:contexto}

El cáncer de pulmón es la principal causa de mortalidad oncológica a nivel
mundial. Según las estimaciones de la \emph{International Agency for Research
on Cancer} (IARC), en 2022 se diagnosticaron aproximadamente 2,48 millones
de nuevos casos, con 1,82 millones de fallecimientos atribuibles a la
enfermedad~\cite{sung2022global}. En España, el Observatorio del Cáncer de la
Asociación Española Contra el Cáncer sitúa la incidencia en torno a 30\,000
casos anuales, con una supervivencia a cinco años que apenas alcanza el
20\,\%. Estas cifras contrastan con las tasas de supervivencia de otros
tumores sólidos como el cáncer de mama, próstata o colorrectal, y evidencian
la necesidad de mejorar tanto el diagnóstico temprano como la
estratificación molecular de los pacientes.

Desde el punto de vista histológico, el cáncer de pulmón se clasifica
principalmente en dos grandes grupos: el cáncer de pulmón de células no
pequeñas (\emph{Non-Small Cell Lung Cancer}, NSCLC), que representa
aproximadamente el 85\,\% de los casos, y el cáncer de pulmón microcítico
(\emph{Small Cell Lung Cancer}, SCLC), con el 15\,\% restante. Dentro del
NSCLC, los dos subtipos histológicos predominantes son el adenocarcinoma
(ADC) y el carcinoma escamoso (\emph{Squamous Cell Carcinoma}, SQC),
seguidos a distancia por el carcinoma de células grandes y los tumores
neuroendocrinos de gran célula.

La distinción entre adenocarcinoma y carcinoma escamoso tiene una
implicación terapéutica directa que trasciende el mero interés
clasificatorio. El pemetrexed, un antifolato aprobado para el tratamiento
del NSCLC no escamoso, presenta una eficacia significativamente menor en
histología escamosa y su uso está desaconsejado en este subtipo. El
bevacizumab, un anticuerpo monoclonal dirigido contra el factor de
crecimiento del endotelio vascular (VEGF), está formalmente
contraindicado en carcinoma escamoso por el riesgo de hemorragia pulmonar
grave. Diversos regímenes de inmunoterapia con inhibidores de \emph{immune
checkpoints} (PD-1/PD-L1) muestran también patrones de respuesta
diferentes según el subtipo. Por tanto, una clasificación histológica
inequívoca condiciona la decisión terapéutica y el pronóstico.

En la práctica clínica, la clasificación se realiza sobre biopsia mediante
inmunohistoquímica (IHC), usando paneles de marcadores canónicos:
citoqueratinas y factores de transcripción de linaje epitelial escamoso
(KRT5, TP63, DSG3\ldots) para el diagnóstico de carcinoma escamoso, y
marcadores del programa alveolar tipo II (napsina A, TTF-1/NKX2-1,
surfactantes\ldots) para el adenocarcinoma. Sin embargo, cuando la biopsia
es escasa o el patrón inmunohistoquímico es ambiguo, la decisión puede
demorarse o quedar sujeta a discrepancia entre patólogos. Un panel
molecular objetivo basado en la expresión de un número reducido de genes
podría complementar la IHC en estos casos.

\section{Transcriptómica y bases de datos públicas}
\label{sec:transcriptomica}

La transcriptómica es la disciplina que estudia el conjunto de tránscritos
de ARN producidos por un genoma en un momento y contexto biológico
concretos. A diferencia del genoma, que es esencialmente estable a lo largo
de la vida del individuo, el transcriptoma refleja el estado funcional de
la célula: qué genes se están expresando, en qué proporción y bajo qué
programa regulatorio. Esta característica hace que el transcriptoma sea un
sustrato especialmente informativo para el estudio del cáncer, donde la
transformación maligna se acompaña de reprogramaciones profundas de la
expresión génica~\cite{hanahan2011hallmarks}.

Las dos tecnologías de referencia para cuantificar el transcriptoma son los
\emph{microarrays} de expresión (predominantes hasta aproximadamente
2015) y la secuenciación masiva de ARN (RNA-Seq, dominante desde entonces).
En \emph{microarrays}, un chip fija miles de sondas de oligonucleótidos
complementarias a cada gen; el ARN de la muestra se marca con
fluorescencia y se hibrida al chip, obteniendo una medida de la abundancia
relativa de cada gen. En RNA-Seq, el ARN se convierte a ADN complementario,
se secuencia y las lecturas se mapean al genoma o transcriptoma de
referencia. Ambas tecnologías han generado un enorme volumen de datos
públicos.

El repositorio NCBI Gene Expression Omnibus (GEO,
\url{https://www.ncbi.nlm.nih.gov/geo/}) es el mayor archivo público de
datos de expresión génica~\cite{barrett2013geo}. Cuenta con más de un
millón de muestras de miles de estudios cubriendo prácticamente cualquier
tejido, condición o especie. Para cáncer de pulmón se dispone de centenares
de estudios que abarcan tanto adenocarcinoma como escamoso, tanto muestras
tumorales como controles adyacentes, y tanto plataformas de
\emph{microarrays} clásicas (Affymetrix HG-U133 Plus 2.0, plataforma GPL570)
como datos de RNA-Seq de The Cancer Genome Atlas (TCGA).

Sin embargo, la disponibilidad no equivale a usabilidad. Un investigador que
quiera integrar varias cohortes debe resolver al menos cuatro problemas
técnicos:

\begin{enumerate}
  \item \textbf{Heterogeneidad de plataformas.} Las sondas de un chip
  Affymetrix HG-U133 Plus 2.0 no son directamente comparables con las de un
  chip Illumina o con las cuentas normalizadas de RNA-Seq. Es preciso
  mapear cada sonda a un símbolo génico y agregar sondas múltiples cuando
  varias apuntan al mismo gen.
  \item \textbf{Curación clínica.} Los metadatos de GEO están redactados en
  texto libre por el laboratorio que deposita el estudio. Una misma
  categoría clínica puede aparecer como \emph{lung adenocarcinoma},
  \emph{ADC}, \emph{non-small cell lung cancer, adenocarcinoma histology},
  \emph{primary tumor: adenocarcinoma}, etc. El grupo control puede
  denominarse \emph{normal}, \emph{healthy}, \emph{non-tumor adjacent
  tissue}, \emph{control lung}, o incluso no estar etiquetado
  explícitamente en el campo esperado.
  \item \textbf{Efecto lote.} La variabilidad técnica entre plataformas,
  laboratorios y protocolos introduce fuentes de varianza que pueden
  dominar sobre la señal biológica y confundir cualquier análisis
  posterior.
  \item \textbf{Validación externa.} Una firma génica identificada en una
  cohorte concreta rara vez se replica en cohortes independientes. La
  literatura acumula cientos de firmas propuestas que no han conseguido
  transferirse a la clínica~\cite{ioannidis2005false}.
\end{enumerate}

\section{Modelos de lenguaje aplicados a curación de metadatos}
\label{sec:llm}

Los recientes modelos de lenguaje de gran tamaño (LLMs), como GPT-4, Claude
o Llama, han demostrado una capacidad notable para tareas de comprensión de
texto en dominios especializados, incluyendo la biomedicina. La capacidad
de estos modelos para interpretar texto libre, extraer entidades y
clasificarlas según categorías predefinidas los convierte en herramientas
naturales para automatizar tareas de curación que hasta ahora requerían
intervención manual.

En el ámbito de este trabajo, un LLM puede leer la descripción textual de
una muestra de GEO (por ejemplo, \emph{lung adenocarcinoma, stage IIB,
72-year-old male, former smoker}) y devolver una etiqueta canónica del
grupo experimental (\texttt{enfermo}, \texttt{sano}, \texttt{ambiguo}). Esta
delegación permite integrar cohortes con nomenclaturas heterogéneas sin
tener que codificar reglas específicas para cada estudio, y hace escalable
la incorporación de nuevas cohortes al análisis. El presente trabajo emplea
Llama~3.3-70b servido vía la infraestructura de inferencia de Groq, que
ofrece una latencia baja y un coste marginal accesible para el volumen de
muestras del proyecto.

\section{Alcance y planteamiento del trabajo}
\label{sec:alcance}

Este Trabajo Fin de Máster desarrolla \framework, un \emph{framework} de
análisis transcriptómico y aprendizaje automático que integra los
componentes descritos en las secciones anteriores para responder a una
pregunta concreta: \emph{¿es posible identificar, mediante integración
automatizada de cohortes públicas y validación externa, una firma génica de
biomarcadores replicable entre cohortes independientes para
adenocarcinoma y carcinoma escamoso de pulmón?}

El planteamiento del trabajo se articula sobre tres pilares:

\begin{itemize}
  \item \textbf{Automatización de la integración multi-cohorte} mediante un
  modelo de lenguaje que resuelve la curación clínica sin intervención
  humana, escalando la incorporación de nuevas cohortes.
  \item \textbf{Validación externa exigente} basada en Leave-One-Dataset-Out
  (LODO), donde cada cohorte se retira una vez del entrenamiento y se
  utiliza como test independiente, y en meta-análisis por consenso: un gen
  entra en la firma final sólo si mantiene el mismo signo de cambio en las
  tres cohortes evaluables simultáneamente.
  \item \textbf{Interpretabilidad biológica} de la firma resultante,
  contrastada con marcadores clínicos ya establecidos en
  inmunohistoquímica diagnóstica y con ejes biológicos conocidos del
  cáncer de pulmón (linaje escamoso frente a alveolar, actividad
  proliferativa).
\end{itemize}

El resto de la memoria se organiza como sigue. El capítulo~\ref{cap:estado}
revisa el estado del arte en integración multi-cohorte, firmas génicas
publicadas y aplicaciones de LLMs en bioinformática. El
capítulo~\ref{cap:objetivos} formula los objetivos general y específicos del
trabajo. El capítulo~\ref{cap:metodos} detalla los materiales y métodos, con
particular atención al diseño del \emph{pipeline}. El
capítulo~\ref{cap:resultados} presenta los resultados de las dos tareas de
clasificación, la firma génica validada y los ejes biológicos identificados.
El capítulo~\ref{cap:discusion} discute las implicaciones clínicas y
metodológicas, sitúa los resultados en el contexto de la literatura y
enumera las limitaciones del trabajo. El capítulo~\ref{cap:conclusiones}
resume las conclusiones y esboza líneas de trabajo futuro.
